\section{Evaluation Results}\label{sec:eval-results}

We evaluate the three headline analyses of Section~\ref{sec:compute}
against the standard tool for each task: \texttt{deconstruct} versus
\texttt{vg deconstruct}~\cite{garrison2018variation}, \texttt{growth} versus
Panacus~\cite{parmigiani2024panacus}, and \texttt{pav} versus
\texttt{odgi pav}~\cite{guarracino2022odgi}. All runs were performed on a single
node with an AMD Ryzen Threadripper PRO 9955WX (16 physical cores, 32 hardware
threads) and 503\,GiB of RAM. Each command was run at \textbf{1} and \textbf{16}
threads, one at a time, and wrapped in \texttt{/usr/bin/time\,-v} to record
wall-clock time and peak resident set size (RSS). \gfaz reads the compressed
\texttt{.gfaz} container directly; the baselines read the uncompressed GFA
(\texttt{odgi} additionally builds an \texttt{.og} index first, timed
separately).

\begin{table}[t]
\centering
\caption{Evaluation graphs. \gfaz operates on the compressed container, which is
$18$--$35\times$ smaller than the GFA the baselines must read.}
\label{tab:eval-datasets}
\small
\begin{tabular}{lrrr}
\hline
Graph & GFA & \texttt{.gfaz} & Ratio \\
\hline
chr1 (HPRC, smoothed) & 7.58\,GB & 214\,MB & 35.4$\times$ \\
chr6 (HPRC, smoothed) & 4.06\,GB & 115\,MB & 35.4$\times$ \\
\emph{E. coli} (MGC)  & 118\,MB  & 6.4\,MB & 18.4$\times$ \\
\hline
\end{tabular}
\end{table}

The chr1 and chr6 graphs are human pangenome graphs with a linear CHM13
reference, used for \texttt{deconstruct} and \texttt{pav}; the \emph{E. coli}
graph is a 13-strain bacterial pangenome used for \texttt{growth}. (The bacterial
graph is circular and stores only W-line walks, so it is not a meaningful target
for reference-anchored \texttt{deconstruct} nor for \texttt{odgi pav}, which does
not load W-line paths; it is retained as a \texttt{growth} data point.)

\subsection{Deconstruct (GFA\,$\rightarrow$\,VCF)}\label{subsec:eval-deconstruct}

Table~\ref{tab:eval-deconstruct} reports the GFA\,$\rightarrow$\,VCF
deconstruction of the chr1 and chr6 graphs against \texttt{vg deconstruct}.
\gfaz produces the same set of variant sites as \texttt{vg}---the record counts
agree to within $0.13\%$ (chr1: $1{,}593{,}901$ vs.\ $1{,}593{,}958$; chr6:
$1{,}304{,}123$ vs.\ $1{,}302{,}387$), consistent with the
$\sim$$99.99\%$ position concordance established in
Section~\ref{subsec:ce-deconstruct}---while running
$36$--$71\times$ faster and using $3.3$--$3.4\times$ less peak memory. The output
is thread-count invariant: the 1- and 16-thread runs emit byte-identical VCFs.
Because deconstruction must additionally hold the per-site allele sequences in
memory, \gfaz's peak RSS here is comparable to the size of the uncompressed GFA
rather than below it; the lighter analyses in
Sections~\ref{subsec:eval-growth}--\ref{subsec:eval-pav} stay well under it.

\begin{table}[t]
\centering
\caption{\texttt{deconstruct} vs.\ \texttt{vg deconstruct}: wall-clock time and
peak RSS at 1 and 16 threads. \gfaz reads the $35\times$-smaller \texttt{.gfaz};
\texttt{vg} reads the GFA.}
\label{tab:eval-deconstruct}
\small
\setlength{\tabcolsep}{4.5pt}
\begin{tabular}{llrrrr}
\hline
 & & \multicolumn{2}{c}{1 thread} & \multicolumn{2}{c}{16 threads} \\
\cline{3-4}\cline{5-6}
Graph & Tool & Time & RSS & Time & RSS \\
\hline
\multirow{2}{*}{chr1} & \texttt{vg}   & 1245\,s & 30.1\,GB & 805\,s & 30.3\,GB \\
                      & \thiswork     & 34.3\,s & 8.8\,GB  & 11.3\,s & 8.3\,GB \\
\hline
\multirow{2}{*}{chr6} & \texttt{vg}   & 898\,s  & 18.4\,GB & 349\,s & 18.4\,GB \\
                      & \thiswork     & 21.9\,s & 5.7\,GB  & 7.0\,s & 5.2\,GB \\
\hline
\end{tabular}
\end{table}

\subsection{Pangenome Growth}\label{subsec:eval-growth}

\emph{Results pending benchmark completion.}

\subsection{Presence/Absence Variation}\label{subsec:eval-pav}

\emph{Results pending benchmark completion.}
